\section{Related Work}\label{sec:related}
Multi-object tracking has been widely studied, with research focusing on improving object detection, data association, and tracking accuracy under different operating conditions. In aerial videos, these tasks become more challenging because of small objects, camera motion, and frequent changes in scene characteristics. Previous studies have addressed these challenges through improved tracking algorithms, adaptive threshold selection, and dynamic inference techniques. In addition, optimization methods and performance evaluation measures have been used to investigate the trade-offs between tracking accuracy and computational efficiency. This section reviews these research directions and discusses their relevance to the problem addressed in this paper.

\subsection{Tracking-by-Detection Methods in Aerial MOT}
Tracking-by-detection is one of the most widely used approaches in multi-object tracking (MOT). In this framework, objects are detected in individual video frames and then associated across consecutive frames to maintain their identities. Several methods have been developed to improve detection reliability, motion prediction, and data association. However, tracking objects in aerial videos introduces additional challenges that require specialized detection and tracking techniques.

SORT built a compact online pipeline from motion prediction and assignment,
and DeepSORT added appearance features to reduce identity fragmentation
\cite{bewley2016sort,wojke2017deepsort}. ByteTrack first associates
high-confidence detections and then uses low-confidence boxes to recover
objects that would otherwise be lost \cite{zhang2022bytetrack}. Tracktor
reuses the detector's box regression for tracking \cite{bergmann2019tracktor},
and TrackFormer performs detection and association with transformer queries
\cite{meinhardt2022trackformer}. Standard evaluation follows the MOTChallenge
benchmarks \cite{milan2020motchallenge}.

Aerial video adds small targets, changing object counts, camera motion, and
clutter \cite{zhu2022visdrone,du2018uavdt}. UAVMOT addresses the motion of the
UAV camera inside the tracker \cite{liu2022uavmot}, and LSO-YOLO designs a
lightweight real-time detector for small UAV objects \cite{dai2026lsoyolo}.
OATrack \cite{qi2026oatrack} is a recent UAV small-object tracker published in
2026. It weights association by detection quality and uses a YOLO detection
cache. We use OATrack as a second host and keep it unchanged. With the same
detector, AC-MOT improves its results in our pipeline: HOTA rises by 1.325,
MOTA by 6.318, and IDF1 by 2.258 points, and IDS fall by 169
(Sect.~\ref{sec:results}). These works improve the tracker or the detector
itself. AC-MOT asks a different question: which detector setting should be
used for the current scene.

\subsection{Adaptive Threshold Selection in Multi-Object Tracking}

Several studies have investigated adaptive threshold selection to improve multi-object tracking under changing scene conditions. In conventional tracking-by-detection systems, fixed thresholds are commonly used to determine which detections are accepted and how they are associated with existing tracks. However, the reliability of these thresholds may vary when object density, occlusion, and detection confidence change during a video sequence. Adaptive threshold methods attempt to address this limitation by adjusting selected tracking parameters according to the current conditions.

Adaptive ByteTrack adjusts the confidence threshold used within the ByteTrack \cite{ma2023adaptivebytetrack} algorithm to accommodate variations in detection confidence. This can help the tracker handle objects whose detection scores change because of occlusion, small object size, or changes in image quality. Similarly, AdapTrack \cite{shim2024adaptrack} adapts the matching thresholds used during data association. By adjusting these thresholds, the tracker can modify its association decisions according to the available detection and tracking information. Although both approaches introduce flexibility into the tracking process, their adaptation mechanisms operate within the tracker and are designed around particular tracking algorithms.

AC-MOT follows a different approach by adapting the detector operating point rather than the internal parameters of the tracker. The detector operating point consists of three parameters: the detection confidence threshold, the non-maximum suppression (NMS) intersection-over-union (IoU) threshold, and the input resolution. Each parameter influences the detections passed to the tracker. The confidence threshold determines which low-confidence detections are retained, the NMS threshold controls the suppression of overlapping bounding boxes, and the input resolution affects the visibility of small objects and the computational cost of detection. Adjusting these parameters together allows the detector configuration to respond to changes in scene conditions before the tracking stage.

An important distinction is that AC-MOT does not require changes to the tracking algorithm. Instead, it operates as an external control mechanism that selects detector configurations using scene information and previous tracking results. The selected configuration is then applied during detector inference, while the tracker continues to perform its original association and identity management procedures. This design allows AC-MOT to be evaluated with different trackers, including ByteTrack and OATrack, without modifying or retraining either tracker. Therefore, AC-MOT extends adaptive parameter selection beyond tracker-specific thresholds by introducing joint detector operating-point control within the overall MOT pipeline.

\subsection{Adaptive Inference for Object Detection}
Dynamic neural networks change their structure or parameters for each input
\cite{han2022dynamic}, and resolution-adaptive networks spend less computation
on easy inputs \cite{yang2020ranet}. AdaScale selects the input scale of each
video frame with a scale regressor trained on detector features
\cite{chin2019adascale}. ApproxDet switches detection configurations at run
time with accuracy and latency models trained offline \cite{xu2020approxdet}.
Chameleon re-profiles the configuration of video analytics pipelines over time
\cite{jiang2018chameleon}. These adaptive inference methods mainly aim to improve object detection accuracy, reduce computational cost, or achieve a suitable balance between accuracy and processing time. To select appropriate inference configurations, some approaches rely on additional trained models, such as scale predictors or offline accuracy and latency estimators, while others use periodic profiling to respond to changes in video content. However, their adaptation strategies are generally designed for object detection or video analytics rather than the specific requirements of multi-object tracking, where detection errors can also affect object association and identity consistency. In contrast, AC-MOT focuses on improving tracking performance by dynamically adjusting three detector parameters: the confidence threshold, non-maximum suppression (NMS) IoU threshold, and input resolution. Its controller uses lightweight scene information and previous tracking results to select detector configurations without requiring an additional learned model or retraining the existing detector or tracker. Furthermore, AC-MOT is evaluated using tracking-specific metrics, including tracking accuracy, identity preservation, and association quality, with two different tracking algorithms, ByteTrack and OATrack. Table~\ref{tab:related} summarizes the main differences between AC-MOT and the existing adaptive inference approaches.

\begin{table}[t]
\caption{Adaptive methods closest to AC-MOT (BT: ByteTrack).}
\label{tab:related}
\centering\footnotesize
\setlength{\tabcolsep}{3pt}
\begin{tabular}{@{}l>{\raggedright\arraybackslash}p{20mm}>{\raggedright\arraybackslash}p{17mm}>{\raggedright\arraybackslash}p{13mm}@{}}
\toprule
Method & Adapted & Where & Task\\
\midrule
AdaScale \cite{chin2019adascale} & input scale & detector & video detection\\
ApproxDet \cite{xu2020approxdet} & detection configuration & detector & video detection\\
Chameleon \cite{jiang2018chameleon} & pipeline configuration & analytics pipeline & video analytics\\
Adapt.\ BT \cite{ma2023adaptivebytetrack} & confidence threshold & inside BT & MOT\\
AdapTrack \cite{shim2024adaptrack} & matching thresholds & inside the tracker & MOT\\
AC-MOT & confidence, NMS, resolution & detector; tracker unchanged & aerial MOT\\
\bottomrule
\end{tabular}
\end{table}

\subsection{Optimization and Evaluation for Multi-Object Tracking}
Optimization and performance evaluation play an important role in developing reliable multi-object tracking (MOT) systems. The selection of detector and tracking parameters can affect detection accuracy, object association, identity consistency, and computational efficiency. Since improving one aspect of tracking performance may negatively affect another, systematic optimization methods are needed to explore different parameter configurations and identify suitable trade-offs. In addition, evaluating tracking systems requires metrics that measure both detection and identity preservation, together with statistical methods that assess the reliability of observed performance differences.

Optuna provides define-by-run search with the Tree-structured Parzen Estimator
(TPE) sampler \cite{akiba2019optuna}, which was used throughout the
development of AC-MOT. The identity-aware study reports a feasible Pareto
front based on the usual notion of non-dominance \cite{deb2002nsga2}. MOTA
weighs detection and identity errors, IDF1 measures identity precision and
recall, and HOTA balances detection and association \cite{luiten2021hota}.
Paired bootstrap resampling at the sequence level keeps the dependence within
each sequence and gives intervals for system differences
\cite{efron1993bootstrap}.
